% Part: first-order-logic
% Chapter: lindstrom
% Section: abstract-logics

\documentclass[../../../include/open-logic-section]{subfiles}

\begin{document}

\olfileid{mod}{lin}{alg}
\olsection{抽象逻辑}

\begin{defn}
一个\emph{抽象逻辑}是一个序对 $\tuple{L, \models_L}$，其中 $L$
是一个函数，它给每个语言~$\Lang{L}$ 指派一个语句集
$L(\Lang{L})$，而 $\models_L$ 是 $\Lang{L}$ 的结构与
$L(\Lang{L})$ 中元素间的关系。特别地，$\tuple{F, \models}$ 就是通常的
一阶逻辑，也就是说，$F$ 是将语言~$\Lang{L}$ 映射到由 $\Lang{L}$ 中常项构成的一阶语句集合的函数，而 $\models$ 是一阶逻辑的满足关系。
\end{defn}

注意，对给定的语言，我们仍采用与一阶逻辑相同的结构概念，
但我们并不预设语句是由 $\Lang{L}$ 中的基本符号按通常方式构建的，
也不预设关系 $\models_L$ 是以与一阶逻辑相同的方式递归定义的。
因此，举例来说，这个完全一般的定义旨在涵盖以下情形：
$\tuple{L,\models_L}$ 中的语句包含无限长的合取或析取，
或者包含除 $\lexists$ 和 $\lforall$ 以外的量词（例如，“存在无穷多个 $x$ 使得……”），
或者可能包含无限长的量词前缀。
为了强调 $L(\Lang{L})$ 中的“语句”不必是通常的一阶逻辑语句，
本章我们使用变元 $!E$、$!F$、~\dots 来遍历它们，
而将 $!A$、$!B$、~\dots 保留给通常的一阶公式。

\begin{defn}
令 $\Mod(L){!E}$ 表示类 $\Setabs{\Struct{M}}{\Struct{M}
  \models_L !E}$。如果需要显式标明语言，我们记作 $\Mod[L](L){!E}$。
两个 $\Lang{L}$-结构 $\Struct{M}$ 和 $\Struct{N}$ 在 $\tuple{L, \models_L}$ 中是\emph{初等等价的}，记作 $\Struct{M} \elemequiv[L] \Struct{N}$，如果它们满足 $L(\Lang{L})$ 中相同的语句。
\end{defn}

\begin{defn}
语言 $\Lang{L}$ 上的抽象逻辑 $\tuple{L,\models_L}$ 被称为\emph{正规的}，如果它满足下列性质：
\begin{enumerate}
\item (\emph{$L$-单调性}) 对语言 $\Lang{L}$ 和 $\Lang{L'}$，若 $\Lang{L} \subseteq \Lang{L'}$，则 $L(\Lang{L}) \subseteq L(\Lang{L'})$。
\item (\emph{扩张性质}) 对每个 $!E \in L(\Lang{L})$，存在 $\Lang{L}$ 的一个\emph{有限}子集 $\Lang{L'}$，使得关系 $\Struct{M} \models_L !E$ 仅依赖于 $\Struct{M}$ 在 $\Lang{L'}$ 上的归约；换言之，若 $\Struct{M}$ 和 $\Struct{N}$ 在 $\Lang{L'}$ 上的归约相同，则 $\Struct{M} \models_L !E$ 当且仅当 $\Struct{N} \models_L !E$。
\item (\emph{同构性质}) 若 $\Struct{M} \models_L !E$ 且 $\Struct{M} \simeq \Struct{N}$，则亦有 $\Struct{N} \models_L !E$。
\item (\emph{改名性质}) 关系 $\models_L$ 在符号改名下保持：若语言 $\Lang{L}'$ 是由 $\Lang{L}$ 将每个符号 $P$ 替换为同元数的符号 $P'$，且每个常项 $c$ 替换为互不相同的常项 $c'$ 而得到的，则对每个结构~$\Struct{M}$ 和语句~$!E$，$\Struct{M} \models_L !E$ 当且仅当 $\Struct{M}' \models_L !E'$，其中 $\Struct{M}'$ 是与 $\Lang{L}$ 对应的 $\Lang{L}'$-结构，且 $!E' \in L(\Lang{L}')$。
\item (\emph{布尔性质}) 抽象逻辑 $\tuple{L, \models_L}$ 对布尔联结词封闭，意即：对每个 $!E \in L(\Lang{L})$，存在 $!F \in L(\Lang{L})$ 使得 $\Struct{M} \models_L !F$ 当且仅当 $\Struct{M} \not\models_L !E$；并且对任意 $!E$ 和 $!F$，存在 $!G$ 使得 $\Mod(L){!G} = \Mod(L){!E} \cap
  \Mod(L){!F}$。对于原子公式和其他联结词也有类似的结论。
\item (\emph{量词性质}) 对 $\Lang{L}$ 中的每个常项 $c$ 和每个 $!E \in L(\Lang{L})$，存在 $!F \in L(\Lang{L})$ 使得
  \[
  \Mod[L'](L){!F} = \Setabs{\Struct{M}}{\Expan{M}{a} \in
  \Mod[L](L){!E} \text{ 对某个 } a \in \Domain{M}},
  \]
  其中 $\Lang{L'} = \Lang{L} \setminus \{c\}$，且 $\Expan{M}{a}$ 是将 $a$ 指派给 $c$ 而得到的 $\Struct{M}$ 到 $\Lang{L}$ 的扩张。
\item (\emph{相对化性质}) 给定一个语句 $!E \in L(\Lang{L})$ 以及不属于 $\Lang{L}$ 的符号 $R$、$c_1$、\dots、$c_n$，存在一个语句 $!F \in L(\Lang{L} \cup \{R,c_1,\ldots,c_n\})$，
  称为 $!E$ 对 $\Atom{R}{x, c_1, \dots c_n}$ 的\emph{相对化}，使得对每个结构~$\Struct{M}$：
  \[
  \Expan{M}{X, b_1, \ldots, b_n} \models_L !F \text{ 当且仅当 } \Struct{N} \models_L !E,
  \]
  其中 $\Struct{N}$ 是 $\Struct{M}$ 的以 $\Domain{N} = \Setabs{a\in \Domain{M}}{\Assign{R}{M}(a, b_1, \dots, b_n)}$ 为论域的子结构
  （参见 \olref[bas][sub]{rem:substructure}），且 $\Expan{M}{X, b_1, \ldots,
    b_n}$ 是将 $\Struct{M}$ 中 $R$， $c_1$，\dots，$c_n$ 分别解释为 $X$， $b_1$，\dots，$b_n$ 而得到的扩张（其中 $X \subseteq M^{n+1}$）。
\end{enumerate}
\end{defn}

\begin{defn}
给定两个抽象逻辑 $\tuple{L_1, \models_{L_1}}$ 与
$\tuple{L_2, \models_{L_2}}$，若对每个语言 $\Lang{L}$
及语句 $!E \in L_1(\Lang{L})$，都存在语句 $!F
\in L_2(\Lang{L})$ 使得 $\Mod[L](L_1){!E} =
\Mod[L](L_2){!F}$，则称后者的表达力至少与前者一样强，记作 $\tuple{L_1, \models_{L_1}}
\leq \tuple{L_2, \models_{L_2}}$。
若 $\tuple{L_1,
  \models_{L_1}} \leq \tuple{L_2, \models_{L_2}}$ 且 $\tuple{L_2,
  \models_{L_2}} \leq \tuple{L_1, \models_{L_1}}$，则称逻辑 $\tuple{L_1, \models_{L_1}}$ 与
$\tuple{L_2, \models_{L_2}}$ 是\emph{等价的}。
\end{defn}

\begin{rem}
  一阶逻辑，即抽象逻辑 $\tuple{F, \models}$，是
  正规的。事实上，上述性质对一阶逻辑来说基本都是直接成立的。我们只须指出，扩张性质可归结为外延性，而将
  语句 $!E$ 相对化到 $\Atom{R}{x, c_1, \dots, c_n}$，是通过
  把其中每个形如 $\lforall[x][!F]$ 的子公式替换为
  $\lforall[x][(\Atom{R}{x, c_1, \dots, c_n} \lif \ !F)]$ 来实现的。此外，
  若 $\tuple{L, \models_L}$ 是正规的，则
  $\tuple{F, \models} \leq \tuple{L, \models_{L}}$，这可以通过对一阶公式施归纳来证明。
  因此，不失一般性，我们可以假设每个一阶
  语句都属于每一个正规逻辑。
\end{rem}

\end{document}